\documentclass[11pt,a4paper]{article}

\usepackage[T1]{fontenc}
\usepackage{lmodern}
\usepackage[margin=28mm]{geometry}
\usepackage{amsmath,amssymb,amsthm}
\usepackage{booktabs}
\usepackage{microtype}
\usepackage[hidelinks]{hyperref}

\newtheorem{theorem}{Theorem}
\newtheorem{lemma}[theorem]{Lemma}
\newtheorem{corollary}[theorem]{Corollary}
\newtheorem{proposition}[theorem]{Proposition}
\newtheorem{problem}[theorem]{Problem}
\theoremstyle{definition}
\newtheorem{definition}[theorem]{Definition}
\theoremstyle{remark}
\newtheorem{remark}[theorem]{Remark}
\newcommand{\Ostar}{O^{*}}
\newcommand{\E}{\mathbb{E}}
\newcommand{\sgn}{\operatorname{sgn}}
\newcommand{\dS}{\bar d_S}

\clubpenalty=10000
\widowpenalty=10000
\displaywidowpenalty=10000
\emergencystretch=2em
\hypersetup{pdftitle={Max 2-SAT: Biased Formulas, the Race for Every Branching Rule, and the Adaptive Direction}}

\title{Max 2-SAT: Biased Formulas, the Race for Every Branching Rule,\\ and the Adaptive Direction}
\author{}
\date{8 October 2026\\[3pt]\small Working notes}

\begin{document}
\maketitle

\begin{abstract}
We take stock of the notes written since 22 September, keep one framework,
and prove one new theorem inside it. The framework is the search with pruning
clauses and the sum check, counted by Knuth's identity and analysed through the
race identity $Z=S+Q-U$. We show that this analysis does not need a fixed
order: along the uniform walk of any branching rule, adaptive or not, the gap
of every variable performs a fair $\pm1$ walk, the shortfall is a fair-coin
thinning and $Q$ is a martingale. The theorem is for \emph{biased} formulas:
if a linear fraction of the variables have an effective gap at least $\varepsilon$
times their own degree, and the degrees have bounded spread, then Max 2-SAT is
solved in time $\Ostar(2^{(1-\gamma)n})$ and polynomial space with
$\gamma=\varepsilon^2\eta^2/(3000R)$, at every clause density, by every
branching rule. No cancellation and no spectral condition is used: a gap
potential $\Phi=\sum_v|\Delta_v|$ is a submartingale, and bias beats the
uncertainty of the last variables. This settles, for a different exponent and
without the measure, the accumulation question left open in the note of
6--7~October. We then organize the remaining work into three regimes, biased,
balanced and aligned, each with its own tool, and describe the adaptive rule
that the race suggests: branch on the clause with the largest shortfall whose
bad values disagree most with the current aggregate of bad values. The slack of
the summed inequality and its exact completion count are proposed as the
measure; they generalize the clause weights $r-\log_\lambda V(r,k)$ to
overlapping clauses with an exact Pascal recursion.
\end{abstract}

\section{Where we stand}\label{sec:state}

\paragraph{Goal.} A deterministic polynomial-space algorithm for Max 2-SAT in
time $\Ostar(c^n)$, $c<2$, for formulas of every clause density; failing that,
natural classes of formulas of unbounded density on which such a bound is
proved for every member.

\paragraph{What is proved.} Table~\ref{tab:status} lists the pieces by the
note that contains them. The rules are all sound for every branching order:
domination, pruning clauses with cuts and propagation~\cite{N3}, sums and
weighted sums of pruning clauses~\cite{N4}, chained pair decisions~\cite{N4},
the three-state compression of near-twin pairs~\cite{NT}, and counting over
exact interaction classes~\cite{N5}. The running-time results are of four
kinds.

\begin{table}[ht]
\centering\small
\begin{tabular}{@{}p{0.30\textwidth}p{0.36\textwidth}p{0.28\textwidth}@{}}
\toprule
Result & Class & Tool \\
\midrule
$2^{H((1-\varepsilon)/2)n}$, tight~\cite{N2,N3} & $\eta n$ variables with gap $\ge\varepsilon n$ & level budget; order-free with pruning clauses\\
Key Lemma, $\gamma<0.0236$~\cite{N7,N8} & random signs on any graph, any density & race identity, late signs, Knuth\\
spectral class~\cite{N9} & $\|J\|,\|L\|\le K\sqrt d$, $\sum h^2\le Knd$ & same race, Talagrand \\
early enumeration~\cite{EE} & random signs, degree spread $D/d$ & count completions of the summed inequality\\
counting criterion~\cite{CC,CA} & (C1) pair counts, (C2) gaps, (C3$'$) & same, deterministic\\
near twins~\cite{NT} & $\beta n$ disjoint $r$-near-twin pairs & compression, $2^n(3/4)^{\beta n/(r+1)}$\\
exact classes~\cite{N5} & $t$ interaction classes, exceptions $E$ & counting ones, $2^{|E|}(n+1)^t$\\
\textbf{this note, Thm.~\ref{thm:biased}} & $\eta n$ variables with gap $\ge\varepsilon D_v$, degree spread $R$ & race for every rule, gap potential\\
\bottomrule
\end{tabular}
\caption{Running-time results. The first row needs gaps linear in $n$; the
new row needs gaps linear in the variable's own degree, so it covers every
density including bounded degree.}\label{tab:status}
\end{table}

\paragraph{What is open.} The measure-and-conquer analysis with clause
weights~\cite[Sections~6--8]{N4} is tight step by step but needs plain steps to
be rare along the tilted walk; the accumulation over many clauses
(\cite[Problem~17]{N4}) is open. The balanced regime needs cancellation of bad
values at the free variables, which random signs provide and which the pair
counts of~\cite{CC} certify; no clause-level condition without a quadratic form
is known. The aligned regime (Max-Cut on $K_m$, the circulant $C_{n,c}$) defeats
every rule implied by single flips~\cite[Section~9]{N4}; near-twin compression
and exact classes solve the extreme cases, and the inverse question, whether a
failure of cancellation forces such structure, is open~\cite{NT}.

\paragraph{What to park.} The learning dichotomy~\cite{N6}: packing and core
certificates, exposure, the entropy dichotomy. Its two lasting points are the
$K_m$ barrier (also in~\cite{N4}) and the joint refutation test, which is the
linear-relaxation cut of~\cite[Proposition~7]{N4}. From the frontier
draft~\cite{N5} we keep the full-phase weights, the weighted-count certificate
and exact interaction classes; the coverage dichotomy is subsumed by the three
regimes below. The spectral class~\cite{N9} is superseded by~\cite{CC,CA} and,
for the biased regime, by Theorem~\ref{thm:biased}, which needs no norm at all.

\section{The race for every branching rule}\label{sec:race}

\paragraph{Setting.} The formula is canonical with $B=1$: unit clauses with
multiplicity, at most one binary clause per pair~\cite[Section~2]{N3}. For a
variable $v$ let $D_v$ be its binary degree, $u_v$ its unit multiplicity after
cancelling opposite units, and $h_v=\Delta_v(F)$ its effective gap at the root:
positive minus negative binary occurrences plus twice the unit imbalance. Then
$|h_v|\le D_v+2u_v$. Write $\bar d=\frac1n\sum_vD_v$ and, for a set $S$,
$\dS=\frac1{|S|}\sum_{v\in S}D_v$. For $a\in\{\pm1\}^n$ the score is
$c_0+\frac14\sum_vh_va_v-\frac14\sum_{u<v}J_{uv}a_ua_v$ with
$J_{uv}\in\{0,\pm1\}$~\cite[Definition~1]{N9}.

\paragraph{The search.} At a node some free variable $x$ is chosen by a
\emph{branching rule}, a deterministic function of the node; both values are
tried. Let $\Delta$ be the effective gap of $x$ at the node, $s=\sgn\Delta$ with
$\sgn0=+1$ its preferred value in $\pm1$ notation, and $p=(D^f-|\Delta|)/2$ the
smaller deficit, where $D^f$ is the number of free neighbours. The child with
$a_x=-s$ receives the pruning clause $\langle L_x\le p-1\rangle$ over the free
neighbours of $x$~\cite[Section~8]{N3}; the other child receives nothing. A
node is cut when a clause is violated (this includes $p\le0$) or when the sum
check fires: $\sum_{w\text{ free}}\min\{m^0_w,m^1_w\}>\sum_ck_c$ over the live
non-vacuous clauses, with $k_c$ the remaining allowance of $c$ and $m^b_w$ the
number of such clauses in which $b$ is the bad value of $w$~\cite[Section~5]{N4}.
No domination and no propagation. Soundness for every rule: each branching
adds a clause to one child only~\cite[Proposition~20]{N3}, and the sum is
implied by the clauses.

\begin{proposition}[Knuth's identity, any rule]\label{prop:knuth}
Let the walk choose a uniform child at every live node, and let $\pi_t$ be the
probability that it is alive after $t$ steps. The tree has $2^t\pi_t$ live
nodes at depth $t$ and at most $1+2n\max_t2^t\pi_t$ nodes in all. Hence the
search takes time $\Ostar(\max_t2^t\pi_t)$ and polynomial space.
\end{proposition}

\begin{proof}
Every live node has two children and is reached by the walk with probability
$2^{-t}$, whatever the rule.
\end{proof}

Write $x_t$ for the variable chosen at step $t$ of the walk, $\sigma_t$ for its
value, $a_t=2\sigma_t-1$, $\Delta_t$ for its gap at that node and
$\xi_t=[a_t=-\sgn\Delta_t]$ for the event that a clause is created. Let
$H_{t}$ be the history $(x_1,\sigma_1,\dots,x_t,\sigma_t)$. The \emph{uncut
process} runs the walk through all $n$ steps, creates every clause and keeps
it. On the event that the walk is alive at depth $T$, the uncut process and
the search agree up to depth~$T$.

\begin{lemma}[The race]\label{lem:race}
At any node, over all clauses created so far,
\begin{gather*}
 \sum_{w\text{ free}}\min\{m^0_w,m^1_w\}-\sum_ck_c=S+Q-U,\\
 S=\sum_c\Bigl(\frac{|\Delta_c|}2+1\Bigr),\qquad
 Q=\sum_c\Bigl(j_c-\frac{a_c}2\Bigr),\qquad
 U=\frac12\sum_{w\text{ free}}|m^1_w-m^0_w|,
\end{gather*}
where $\Delta_c$ is the gap at the creation of $c$, $a_c$ the number of its
support variables assigned since, and $j_c$ the number of those that took their
bad value. If the walk is alive at depth $T$, then $Z_T=S_T+Q_T-U_T\le0$.
\end{lemma}

\begin{proof}
This is~\cite[Lemma~2, Remark~3]{N7}: the identity is algebra at one node and
does not mention the order; removing vacuous clauses only strengthens the
check, so the check of the search fires whenever the check over all created
clauses fires.
\end{proof}

\begin{lemma}[The values are the coins]\label{lem:coins}
For every branching rule:
\begin{enumerate}
 \item[(a)] Given $H_{t-1}$, the variable $x_t$, its gap $\Delta_t$, the
 counts $m^b_{x_t}$ and the free set are determined, and $a_t$ is uniform on
 $\pm1$. Hence $\xi_t$ is a fair coin given $H_{t-1}$.
 \item[(b)] The gap of a free variable $v$ is
 $\Delta_v(t)=h_v-\sum_{s\le t,\,x_s\sim v}J_{vx_s}a_s$. When a neighbour
 $x_s$ of $v$ is assigned, $\Delta_v$ changes by $\pm1$, each with probability
 $\frac12$ given $H_{s-1}$. So, sampled at the times its neighbours are
 assigned, $\Delta_v$ is $h_v$ plus a simple symmetric random walk.
 \item[(c)] $Q_t-Q_{t-1}=\frac12a_tb_{x_t}$ with $b_{x_t}=m^1_{x_t}-m^0_{x_t}$
 taken at depth $t-1$; so $Q$ is a martingale with
 $|Q_t-Q_{t-1}|\le\frac12D_{x_t}$.
\end{enumerate}
\end{lemma}

\begin{proof}
(a) The rule is a function of the node, and $\sigma_t$ is a fresh uniform bit.
(b) The coefficient of $a_v$ in the score after substituting the assigned
variables is $\frac14(h_v-\sum J_{vx_s}a_s)$, and in the residual canonical
formula it is $\frac14\Delta_v$~\cite[Lemma~4(a)]{N9}; the increment
$-J_{vx_s}a_s$ is $\pm1$ with probability $\frac12$ because $a_s$ is uniform
given $H_{s-1}$ and $J_{vx_s}=\pm1$ is fixed. A $\pm1$ sequence in which every
term is uniform given the past is a sequence of independent fair coins.
(c) Assigning $x_t$ raises $a_c$ by one for every created clause containing
it and $j_c$ by $[\sigma_t\text{ bad in }c]$, which sums to $\frac12a_tb_{x_t}$;
$b_{x_t}$ is determined by $H_{t-1}$ and $|b_{x_t}|$ is at most the number of
created clauses containing $x_t$, at most $D_{x_t}$.
\end{proof}

This is the deterministic form of the late-signs lemma of~\cite{N8}: there the
unseen coins were the signs, here they are the values. The difference is that
the bad values of the free variables are now fixed by the formula, so $U_T$ is
not a sum of coins. For biased formulas we do not need it to be.

\section{Biased formulas}\label{sec:biased}

\begin{definition}[Biased formulas]\label{def:biased}
Let $\varepsilon,\eta\in(0,1]$ and $R\ge1$. A canonical formula with $B=1$ is
\emph{$(\varepsilon,\eta,R)$-biased} if there is a set $S$ of at least
$\eta n$ variables such that
\begin{enumerate}
 \item[(B1)] every $v\in S$ has $|h_v|\ge\varepsilon D_v$: its majority
 literal exceeds its minority literal by an $\varepsilon$-fraction of its
 degree, counting a unit clause as two occurrences;
 \item[(B2)] $\dS\ge1$ and $\sum_{v\in V}(D_v+u_v)^2\le R\,n\,\dS^{\,2}$: no
 part of the formula is much denser than the biased variables.
\end{enumerate}
\end{definition}

Condition (B1) is the gap condition of~\cite{N4} in its relative form, where
the gap is compared with the variable's own degree and not with $n$; it allows
every density, from bounded degree to complete graphs. Condition (B2) holds
with $R=1$ for regular graphs without unit clauses and with $R=O(1)$ whenever
all degrees are within a constant factor of $\dS$. The set $S$ need not be
known to the algorithm.

\begin{theorem}[Biased formulas, every branching rule]\label{thm:biased}
Let $F$ be $(\varepsilon,\eta,R)$-biased. For every branching rule, the search
of Section~\ref{sec:race} solves Max 2-SAT on $F$ in polynomial space and time
\[
 \Ostar\bigl(2^{(1-\gamma)n}\bigr),\qquad
 \gamma=\frac{\varepsilon^2\eta^2}{3000\,R}.
\]
\end{theorem}

The constant is not optimized; the order $\varepsilon^2$ is right
(Remark~\ref{rem:tight}). The proof uses three estimates along the walk of
Lemma~\ref{lem:coins}, one for each term of the race, and a potential that
turns the initial gaps into a lower bound on the gaps at branching time.

\begin{lemma}[Gap potential]\label{lem:potential}
Fix a set $S$ and a branching rule, and let the uncut process apply the rule
to every node, cut or not. Put $\Phi_t=\sum_{v\in S\text{ free at }t}|\Delta_v(t)|$
and $G_t=\sum_{s\le t,\,x_s\in S}|\Delta_s|$. Then
\[
 G_T=\Phi_0-\Phi_T+\sum_{t\le T}z_t+M_T ,
\]
where $z_t\ge0$ is the number of free neighbours $v\in S$ of $x_t$ with
$\Delta_v(t-1)=0$, and $M$ is a martingale with respect to $(H_t)$ with
$|M_t-M_{t-1}|\le2D_{x_t}$. Consequently
\[
 \Pr\bigl[G_T<\Phi_0-\Phi_T-x\bigr]\le
 \exp\Bigl(-\frac{x^2}{8\sum_{t\le T}D_{x_t}^2}\Bigr).
\]
\end{lemma}

\begin{proof}
At step $t$ the variable $x_t$ leaves the free set, taking $|\Delta_t|$ with it
if $x_t\in S$, and each free neighbour $v\in S$ of $x_t$ changes its gap by
$\pm1$. So $\Phi_t-\Phi_{t-1}=-[x_t\in S]\,|\Delta_t|+Y_t$ with
$Y_t=\sum_v(|\Delta_v(t)|-|\Delta_v(t-1)|)$ over these neighbours. By
Lemma~\ref{lem:coins}(b) each term is $|\Delta\pm1|-|\Delta|$ with a uniform
sign given $H_{t-1}$, whose conditional mean is $[\Delta=0]$; hence
$\E[Y_t\mid H_{t-1}]=z_t$. Summing over $t\le T$ gives
$\Phi_T-\Phi_0=-G_T+\sum_tY_t$, which is the identity with
$M_T=\sum_{t\le T}(Y_t-z_t)$. The number of terms in $Y_t$ is at most
$D_{x_t}$, so $Y_t\in[-D_{x_t},D_{x_t}]$ and $z_t\in[0,D_{x_t}]$, whence
$|Y_t-z_t|\le2D_{x_t}$. The tail is the Azuma--Hoeffding
inequality~\cite{Azu67,Hoe63}.
\end{proof}

The potential says that bias cannot be destroyed: assigning a neighbour moves
a gap by one step of a fair walk, which does not lower $|\Delta_v|$ in
expectation, whatever the rule does. The term $\sum_tz_t$ is the source of the
$\sqrt d$ gaps of balanced formulas (a walk at $0$ moves away); for biased
formulas we drop it.

\begin{lemma}[Shortfall and martingale]\label{lem:tails}
For every branching rule, with $V=\sum_{v}\bigl((D_v+u_v)^2+1\bigr)$ and
$E=\sum_vD_v^2$:
\begin{enumerate}
 \item[(a)] $\Pr\bigl[S_T\le\tfrac14G^{\rm all}_T+\tfrac T2-\alpha\bigr]\le
 e^{-\alpha^2/(2V)}$, where $G^{\rm all}_T=\sum_{t\le T}|\Delta_t|\ge G_T$;
 \item[(b)] $\Pr[Q_T\le-\kappa]\le e^{-2\kappa^2/E}$;
 \item[(c)] $U_T\le\frac12\sum_{w\text{ free at }T}D_w$ for every walk.
\end{enumerate}
\end{lemma}

\begin{proof}
(a) $S_T=\sum_{t\le T}X_t$ with $X_t=\xi_t(|\Delta_t|/2+1)\ge0$. By
Lemma~\ref{lem:coins}(a), $\E[X_t\mid H_{t-1}]=|\Delta_t|/4+\frac12$ and
$\E[X_t^2\mid H_{t-1}]=\frac12(|\Delta_t|/2+1)^2\le(D_{x_t}+u_{x_t})^2+1=:v_t$,
using $|\Delta_t|\le|h_{x_t}|+D_{x_t}\le2(D_{x_t}+u_{x_t})$ and
$(a+1)^2\le2a^2+2$. For $\lambda\ge0$, $e^{-\lambda X}\le1-\lambda X+\lambda^2X^2/2$
gives $\E[e^{-\lambda X_t}\mid H_{t-1}]\le\exp(-\lambda\E[X_t\mid H_{t-1}]
+\lambda^2v_t/2)$, so $\exp(-\lambda\sum_{s\le t}(X_s-\E[X_s\mid H_{s-1}])
-\frac{\lambda^2}2\sum_{s\le t}v_s)$ is a supermartingale, and $\lambda=\alpha/V$
with $\sum_tv_t\le V$ gives the bound~\cite{Mau03}. (b) Lemma~\ref{lem:coins}(c)
and Azuma--Hoeffding with $\sum_t(D_{x_t}/2)^2\le E/4$. (c) $|m^1_w-m^0_w|$ is
at most the number of created clauses containing $w$, at most $D_w$.
\end{proof}

\begin{proof}[Proof of Theorem~\ref{thm:biased}]
Write $d=\dS$ and $\theta=\varepsilon\eta$. Put $\delta=\theta^2/(256R)$ and
$T=\lceil(1-\delta)n\rceil$, so that the free set $W$ at depth $T$ has at most
$\delta n$ variables. By Proposition~\ref{prop:knuth} and
$\max_t2^t\pi_t\le\max\{2^T,2^n\pi_T\}$ (as $\pi_t$ is non-increasing), it
suffices to show $\pi_T\le3e^{-cn}$ with $c=\theta^2/(4096R)$: then
$\gamma=\min\{\delta,c/\ln2\}\ge\theta^2/(3000R)$.

\emph{Sizes.} By (B2), $E=\sum_vD_v^2\le Rnd^2$ and
$V=\sum_v((D_v+u_v)^2+1)\le Rnd^2+n\le2Rnd^2$. By Cauchy--Schwarz, any set
$W$ of at most $\delta n$ variables has
$\sum_{w\in W}(D_w+u_w)\le\sqrt{|W|\sum_v(D_v+u_v)^2}\le\sqrt{\delta R}\,nd
=\tfrac\theta{16}\,nd$.

\emph{The potential.} $\Phi_0=\sum_{v\in S}|h_v|\ge\varepsilon\sum_{v\in S}D_v
=\varepsilon|S|d\ge\theta nd$ by (B1). At depth $T$,
$|\Delta_w(T)|\le|h_w|+D_w\le2(D_w+u_w)$, so
$\Phi_T\le2\sum_{w\in W}(D_w+u_w)\le\tfrac\theta8nd$. Take $x=\tfrac\theta4nd$
in Lemma~\ref{lem:potential}: except with probability
$\exp(-x^2/(8E))\le\exp(-\theta^2n/(128R))$,
\[
 G_T\ \ge\ \theta nd-\tfrac\theta8nd-\tfrac\theta4nd=\tfrac58\,\theta nd .
\]

\emph{The three terms.} Take $\alpha=\kappa=\tfrac\theta{32}nd$ in
Lemma~\ref{lem:tails}. Except with probability
$e^{-\alpha^2/(2V)}\le e^{-\theta^2n/(4096R)}$ we have
$S_T\ge\frac14G_T-\alpha\ge\frac5{32}\theta nd-\frac1{32}\theta nd=\frac4{32}\theta nd$,
and except with probability $e^{-2\kappa^2/E}\le e^{-\theta^2n/(512R)}$ we
have $Q_T\ge-\frac1{32}\theta nd$. For every walk,
$U_T\le\frac12\sum_{w\in W}D_w\le\frac1{32}\theta nd$.

\emph{The race.} Outside the three exceptional events,
$Z_T=S_T+Q_T-U_T\ge\frac{4-1-1}{32}\theta nd>0$, so by Lemma~\ref{lem:race} the
walk is not alive at depth $T$. Hence $\pi_T\le3e^{-\theta^2n/(4096R)}$.
\end{proof}

\begin{remark}[Every density, every rule]
Only (B1) and (B2) enter; $D_v$ may be $3$ or $n-1$. The branching rule is
arbitrary, so the theorem covers in particular the adaptive rule
of~\cite[Section~6]{N4} run without domination and propagation: for formulas
in which every variable has ratio at most $\beta_0<\frac12$, which are
$(1-2\beta_0,1,R)$-biased, the tree has $2^{(1-\gamma)n}$ nodes whatever the
plain steps do. This answers the accumulation question
of~\cite[Problem~17, open question~1]{N4} in the weak form ``some constant
below $2$''; the exponent $H(\beta')$ hoped for there remains open and would
need the measure.
\end{remark}

\begin{remark}[Order of the exponent]\label{rem:tight}
The circulant formula $C_{n,c}$ with gap $\varepsilon n$ is
$(\varepsilon,1,1)$-biased, and every rule of~\cite{N3} has exactly
$\binom{n-1}{a}=2^{H((1-\varepsilon)/2)n-O(\log n)}$ leaves on it; the sum
check does not change this~\cite[Section~5]{N4}. So no analysis of this search
can give $\gamma>1-H((1-\varepsilon)/2)=\varepsilon^2/(2\ln2)+O(\varepsilon^4)$,
and the order $\varepsilon^2$ in Theorem~\ref{thm:biased} is right. The
constant $3000$ comes from three union-bounded tails and the crude bound
$|\Delta_t|\le2(D+u)$; a single exponential moment of $Z_T$ would improve it.
\end{remark}

\begin{remark}[Why the sum check is needed]
For a bounded-degree variable, the probability that it is undominated after
its predecessors are assigned is a constant, so the level budget of~\cite{N3}
cannot certify an exponential saving at the last levels; and a single pruning
clause compares a shortfall $\varepsilon D_x/2$ with its whole free support,
which is decided only at the end. In the race the shortfalls of all clauses add
up, while only the variables still free at depth $T$ carry uncertainty; bias
of order $\varepsilon\bar d$ per clause beats uncertainty of order $\bar d$ per
free variable as soon as fewer than $\varepsilon n/2$ variables are free. This
is the ``accumulation over many clauses'' that the measure with nearly disjoint
counted clauses cannot do~\cite[Remark~12]{N4}.
\end{remark}

\begin{remark}[Better constants from the order]
With the fixed order by decreasing degree and $S=V$, the last $\delta n$
variables have $\sum_{w\in W}D_w\le\delta n\bar d$ without the square root, so
$\delta$ can be taken of order $\varepsilon/R$ and the margin is linear in
$\varepsilon$; only the tails keep the $\varepsilon^2$. With partial bias, put
$S$ last. The theorem is stated for every rule because the adaptive direction
of Section~\ref{sec:adaptive} needs that freedom.
\end{remark}

\begin{remark}[What is not covered]
Domination and propagation are sound but change which values are coins; the
count is for the plain search. Complementary pairs ($B=2$) give steps $\pm2$
and change constants once the race identity with multiplicities is written
out~\cite[Section~4]{N9}. Weighted clauses are not covered.
\end{remark}

\begin{table}[ht]
\centering
\begin{tabular}{llccccc}
\toprule
 & & $d=12$ & $32$ & $128$ & $150$\\
\midrule
random signs, $\varepsilon=\frac14$ & single clauses & 0.86 & 0.77 & 0.57 & 0.60\\
 & with the sum check & 0.90 & 0.91 & 0.89 & 0.93\\
random signs, $\varepsilon=\frac12$ & single clauses & 0.90 & 0.86 & 0.81 & 0.78\\
 & with the sum check & 0.92 & 0.93 & 0.94 & 0.95\\
aligned, $\varepsilon=\frac14$ & single clauses & 0.74 & 0.65 & 0.51 & 0.55\\
 & with the sum check & 0.72 & 0.72 & 0.51 & 0.47\\
aligned, $\varepsilon=\frac12$ & single clauses & 0.84 & 0.73 & 0.68 & 0.70\\
 & with the sum check & 0.85 & 0.76 & 0.76 & 0.69\\
\bottomrule
\end{tabular}
\caption{Saving $n-\log_2(\text{tree size})$ as a fraction of $n$ on biased
formulas with $n=300$ (fixed order, population splitting with $M=40$). Each
pair carries a clause with probability $d/(n-1)$; every variable has
$|h_v|\ge\varepsilon D_v$. ``Aligned'' means that every clause is $(u\lor v)$
or $(\neg u\lor\neg v)$ after renaming each variable to its majority literal,
the type of $C_{n,c}$.}\label{tab:biased}
\end{table}

Table~\ref{tab:biased} shows the two faces of the theorem. With random signs
the sum check keeps the saving flat in $d$ while single clauses decay, exactly
as in the balanced case~\cite[Table~2]{N7}, but at a much higher level. On
aligned formulas the sum has no opposite literals to cancel and gives nothing;
the saving is then the bias alone, and it is still flat in $d$.

\section{Three regimes, one framework}\label{sec:regimes}

The race separates the formulas by what wins it. Write $N^f_x$ for the free
neighbours of $x$ at a node, and for a free variable $w$ let
$b_w=m^1_w-m^0_w$ be the \emph{aggregate}: the excess of live clauses in
which $w=1$ is bad over those in which $w=0$ is bad.

\begin{enumerate}
 \item \textbf{Biased.} Gaps of order $\varepsilon D$ give a shortfall of order
 $\varepsilon\bar d$ per clause, which beats the trivial uncertainty $D_w$ of
 each free variable once few variables are free. Theorem~\ref{thm:biased}:
 every rule, every density, no cancellation.
 \item \textbf{Balanced, unaligned.} Gaps are of order $\sqrt D$ (the term
 $\sum_tz_t$ of Lemma~\ref{lem:potential}: a gap at $0$ moves away), so the
 shortfall is of order $\sqrt d$ per clause and the race is won only if
 $|b_w|$ is of order $\sqrt{m_w}$ rather than $m_w$: the bad values of the
 clauses containing $w$ must cancel. Random signs give this~\cite{N8}; for a
 fixed formula the aggregate pair-count condition (C1) of~\cite{CC} certifies
 it, and Lemma~\ref{lem:potential} replaces the gap-profile condition (C3)
 without Talagrand's inequality and without (C1), under the degree-spread
 condition (B2) alone (Section~\ref{sec:balanced}).
 \item \textbf{Aligned.} Many clauses agree about the bad values of many
 variables, so $b_w$ is of order $m_w$ and the sum is useless; by the barrier
 of~\cite[Section~9]{N4} no rule implied by single flips helps. The exits are
 structural: near-twin compression~\cite{NT} and exact interaction
 classes~\cite{N5}.
\end{enumerate}

The research direction is then a single question.

\begin{problem}[Dichotomy at a node]\label{prob:dichotomy}
Find a polynomial-time test at a node, in terms of gaps, aggregates $b_w$ and
agreement counts between signed neighbourhoods, with the following property.
Either the test names a free variable whose pruning clause raises $Z$ by a
constant fraction of its support size, so that the race is won along the walk;
or it names a set of variables whose signed neighbourhoods agree on all but a
constant fraction of their supports, so that near-twin compression or class
counting solves the residual formula within the budget. Then prove that one of
the two happens at every node of every formula of density $\omega(1)$.
\end{problem}

The near-twin note proves the second exit for pairs and leaves open whether a
failure of cancellation forces linearly many near-twin pairs~\cite[Section~5]{NT}.
The zero-gap Hadamard example of~\cite[Section~6]{CC} shows that pairwise
cancellation and alignment are not the only obstacles: there the bad values
cancel exactly, but no clause has a shortfall. Lemma~\ref{lem:potential}
explains it: with $h=0$ and no earlier neighbours, $\sum_tz_t$ is the only
source of gaps, and the left side, branched first, never has an assigned
neighbour.

\subsection{The balanced regime without Talagrand}\label{sec:balanced}

The deterministic criterion of~\cite{CC} needs (C1) for two things: the
martingale $Q$, the uncertainty $U$, and, through the addendum~\cite{CA}, the
concentration of the total gap. The potential gives the last one directly.

\begin{proposition}\label{prop:gaps-balanced}
For a fixed order, $\E\,G^{\rm all}_T\ge\sum_{t\le T}\max\{|h_t|,\sqrt{D^-_t/3}\}$,
where $D^-_t$ counts the neighbours of $x_t$ before it, and for every $x>0$
\[
 \Pr\bigl[G^{\rm all}_T<\E\,G^{\rm all}_T-x\bigr]\le\exp\Bigl(-\frac{x^2}{8\sum_vD_v^2}\Bigr).
\]
\end{proposition}

\begin{proof}
The mean is the computation of~\cite[Proposition~1]{CA}: $\Delta_t$ is $h_t$
minus a symmetric sum of $D^-_t$ fair signs, so $\E|\Delta_t|\ge|h_t|$ by
Jensen and $\E|\Delta_t|\ge\E|X_t|\ge\sqrt{D^-_t/3}$ by H\"older. For a fixed
order $G^{\rm all}_T$ is a function of the independent values $a_1,\dots,a_T$,
and changing $a_u$ changes the gap of each later neighbour of $x_u$ by $2$, so
$G^{\rm all}_T$ changes by at most $2D_{x_u}$. McDiarmid's inequality with
these differences gives the tail. For an adaptive rule the same bound follows
from Lemma~\ref{lem:potential}, whose martingale has the same increments.
\end{proof}

So under (B2) the condition (C3) of~\cite{CC} follows from (C3$'$) with
$\eta=g^2/(8R)$, in place of $g^2/(128\tau K)$ in~\cite{CA}, and without the
norm bound. What (C1) is still needed for is $Q$, whose steps are
$\frac12b_{x_t}$, and $U_T$: both are sums of bad values, which for a fixed
formula are not coins. In clause terms, (C1) says that for every weighting of
the branched variables the weighted agreement counts between their signed
neighbourhoods stay bounded by $Kd$ times the weight. This is the one place
where a quadratic form has not been replaced by a count of clauses; the
row-sum version fails on random formulas~\cite[Remark~5]{CA}. We regard it as
the definition of ``unaligned'' until Problem~\ref{prob:dichotomy} replaces it
by an adaptive test.

\section{The adaptive rule that the race suggests}\label{sec:adaptive}

Everything the race needs is visible at a node and countable in clause terms.

\subsection{The node}

Let $W$ be the free variables, $\mathcal C$ the live clauses with allowances
$k_c$, and for $w\in W$ let $m^0_w,m^1_w$ count the clauses of $\mathcal C$ in
which $0$, respectively $1$, is the bad value of $w$. The \emph{aggregate} is
$b_w=m^1_w-m^0_w$ and the \emph{slack} is
\[
 B=\sum_{c\in\mathcal C}k_c-\sum_{w\in W}\min\{m^0_w,m^1_w\}=U-S-Q\ \ge0 ,
\]
the node being cut exactly when $B<0$. For a free variable $x$ with gap
$\Delta_x$, preferred value $s=\sgn\Delta_x$ and free neighbours $N^f_x$, the
pruning clause created in the child $a_x=-s$ has bad values
$\beta_x(w)=-sJ_{xw}\in\{\pm1\}$ for $w\in N^f_x$~\cite[Lemma~4(b)]{N9}. Its
\emph{agreement} with the aggregate is
\begin{align*}
 A_x&=\sum_{w\in N^f_x}\bigl(|b_w+\beta_x(w)|-|b_w|\bigr)\\
 &=\#\{w:\beta_x(w)=\sgn b_w\text{ or }b_w=0\}-\#\{w:\beta_x(w)=-\sgn b_w\ne0\},
\end{align*}
and its \emph{value} is $\nu_x=|\Delta_x|/2+1-A_x/2$.

\begin{proposition}[Child slacks]\label{prop:children}
Branching on $x$ gives the children
\[
 B_{\rm pref}=B-\tfrac12\bigl(|b_x|+s\,b_x\bigr),\qquad
 B_{\rm pruned}=B-\tfrac12\bigl(|b_x|-s\,b_x\bigr)-\nu_x .
\]
\end{proposition}

\begin{proof}
In the child $a_x=a$, $U$ loses $\frac12|b_x|$ and $Q$ gains $\frac12ab_x$
(Lemma~\ref{lem:coins}(c)), so $B$ loses $\frac12(|b_x|+ab_x)$. The pruned
child also adds $|\Delta_x|/2+1$ to $S$ and $\frac12A_x$ to $U$.
\end{proof}

So exactly one child, the one in which $x$ takes its majority bad value, pays
$|b_x|$; this is a Pascal step, tight for every weight. The pruned child pays
in addition the value of the new clause, which is positive when the gap is
large and the clause disagrees with the aggregate on its free support, and
negative when the clause agrees with it.

\subsection{The three steps of~\cite{N4} in these terms}

\begin{itemize}
 \item \emph{Case~2, inside the clauses.} The variable with the largest
 $|b_x|$: the most contested free variable, where the live clauses disagree
 most about the bad value. This generalizes ``branch on the clause most likely
 to be violated''~\cite[Section~4]{N4} from one clause to all of them.
 \item \emph{Case~1, a new clause.} The variable with the largest $\nu_x$: a
 large gap, and bad values that disagree with the current aggregate.
 ``Counting only sufficiently new clauses''~\cite[Section~6]{N4} becomes a
 number: a clause is new to the extent that $A_x$ is below its support size,
 and it is worth $\nu_x$ units of slack.
 \item \emph{A plain step.} $b_x=0$ and $\nu_x\le0$ for every free $x$: no
 contested variable and no clause that beats the aggregate.
\end{itemize}

One-step greed, maximizing $|b_x|+\nu_x$, is the wrong currency early in the
tree: a fresh clause has a large free support and $b\approx0$ there, so
$A_x\approx|N^f_x|$ and $\nu_x<0$ for every balanced $x$, although the race is
won in the end. The reason is timing~\cite[Remark~29]{N3}: the support of a
clause is paid into $U$ at creation and repaid through $Q$, as fair coins,
when the support is assigned. The right currency is what is left at the end.

\subsection{A reserve}

Fix in advance a set $W_0$ of $\delta n$ variables, the \emph{reserve}, to be
assigned last. This is the set $L$ of~\cite[Section~7]{N2}, with a weaker
demand: not that every reserved variable be dominated, but that the summed
inequality over $W_0$ have few completions~\cite{EE}. Branch on $V\setminus W_0$
by any rule; at depth $n-\delta n$ enumerate the completions of the summed
inequality over $W_0$, or continue the search inside $W_0$.

\begin{proposition}[What the reserve costs]\label{prop:reserve}
Let the rule on $V\setminus W_0$ be arbitrary. For the uniform walk let
$s=S+Q$ at depth $n-\delta n$, where $Q$ is over the assigned support
variables, and let $E_W=\sum_{w\in W_0}b_w^2$ be the energy of the aggregate
on the reserve. Then the number of completions of the summed inequality over
$W_0$ at that node is at most $2^{\delta n}\rho$, where $\rho=\exp(-2s^2/E_W)$ if
$s>0$ and $\rho=1$ otherwise, and the total work of the search is
$\Ostar\bigl(2^{(1-\delta)n}+2^n\,\E\rho\bigr)$.
\end{proposition}

\begin{proof}
A completion $a\in\{\pm1\}^{W_0}$ satisfies the summed inequality iff
$s+\frac12\sum_{w\in W_0}a_wb_w\le0$; for uniform $a$ this is a Hoeffding
tail~\cite[Lemma~1]{CC}. Knuth's identity gives the total.
\end{proof}

So the rule on $V\setminus W_0$ should make $s^2/E_W$ large: large shortfalls,
and an aggregate on the reserve that stays balanced. Creating the clause of
$x$ changes $s$ by $|\Delta_x|/2+1$ and $E_W$ by
$|N^f_x\cap W_0|+2\sum_{w\in N^f_x\cap W_0}b_w\beta_x(w)$, where the last sum
is the correlation of the new clause with the aggregate on the reserve, the sum
of the pair counts $C_{xc}$ of~\cite[(1)]{CC} over the live clauses $c$.

\paragraph{Rule R (designed cancellation).} Write
$\rho_x=\sum_{w\in N^f_x\cap W_0}b_w\beta_x(w)$ for the correlation of the new
clause with the aggregate on the reserve. To first order, creating the clause
of $x$ changes $s^2/E_W$ by $(s/E_W)\bigl(|\Delta_x|+2-(s/E_W)(|N^f_x\cap W_0|+2\rho_x)\bigr)$.
Among the free variables outside the reserve, branch on the $x$ that
maximizes
\[
 |\Delta_x|+2-\frac{s}{E_W}\bigl(|N^f_x\cap W_0|+2\rho_x\bigr):
\]
a large gap, and bad values that disagree with the aggregate on the reserve.
When every free $x$ has a large positive $\rho_x$, prefer the largest $|b_x|$
instead. On biased formulas any rule wins (Theorem~\ref{thm:biased}); on
balanced formulas the correlation term is the only control over $E_W$ we
have, and it is a count of agreeing minus disagreeing bad values, exactly the
quantity whose aggregate bound is (C1).

\paragraph{The exit.} If at some node every free $x$ outside the reserve has
correlation at least $(1-\kappa)|N^f_x\cap W_0|$ with the aggregate, then every
remaining row $J_{x,W_0}$ agrees with $-s_x\sgn(b)$ on all but a
$\kappa$-fraction of its support: the clauses between the remaining variables
and the reserve are, after renaming the reserve by $\sgn b$, all of the same
type up to exceptions. This is the near-twin structure of~\cite{NT} with
respect to $W_0$; what it does not control are the clauses inside the
remaining set and inside the reserve. Making this exit an algorithm with a
certified cost is Problem~\ref{prob:dichotomy} in its concrete form.

\section{Measures}\label{sec:measures}

The analyses above use five quantities as measures. None of them is additive
over variables, and that is the point: additive measures certify a saving
branch by branch, which fails for linear gaps~\cite[Section~7.2]{N3}.

\begin{enumerate}
 \item \textbf{Knuth's measure} $\max_t\,(t+\log_2\pi_t)$. It is exact for
 every search in which every live node has two children, it absorbs plain
 steps, and it reduces the analysis to one survival probability.
 Theorem~\ref{thm:biased} is an analysis of this measure; so is the Key
 Lemma~\cite{N8}.
 \item \textbf{The volume of the summed inequality.} At a node orient every
 free $w$ so that its cheaper value is $0$, let $c_w=|b_w|$ and
 $V(c,B)=\#\{z\in\{0,1\}^W:\sum_wc_wz_w\le B\}$, computable by the knapsack
 dynamic program in polynomial time~\cite[Section~3]{N5}. Put
 \[
  \mu_\lambda=\#\{w\in W:b_w=0\}+\log_\lambda V(c,B),
 \]
 so that $\mu_\lambda=n$ at the root. For $c_w=1$ on $r$ variables and $B=k$
 this is $r-e_\lambda(r,k)$ with the weights $e_\lambda(r,k)=r-\log_\lambda V(r,k)$
 of~\cite[Theorem~10]{N4}; the measure extends them to overlapping clauses by
 summing them first. A step on a contested variable $w$ satisfies
 $V(c,B)=V(c\setminus w,B)+V(c\setminus w,B-c_w)$, so it is tight for every
 $\lambda$ (Pascal's rule); a step on an uncontested variable has drops $1$
 and $1$; a creation is admissible at base $\lambda$ iff it raises
 $e_\lambda=|W|-\log_\lambda V$ by at least $\log_\lambda\frac1{\lambda-1}$,
 which is the condition of~\cite[Lemma~7]{N4} with the sum in place of one
 clause. So ``a clause counts if it is sufficiently new'' has an exact
 meaning: it counts if it lowers the volume of the summed inequality by a
 factor $\lambda^{-c(\lambda)}$. With $\lambda=2$ this is the certificate
 $N_{\mathbf 1}$ of~\cite{N5} and the quantity that the reserve enumeration
 pays~\cite{EE}; with $\lambda<2$ it is the measure of~\cite{N4}, and the
 plain steps are again the obstacle.
 \item \textbf{The gap potential} $\Phi=\sum_{v\text{ free}}|\Delta_v|$, a
 submartingale along the walk of every rule (Lemma~\ref{lem:potential}). It
 certifies future shortfall: the bias in the formula is conserved in
 expectation and is spent, $|\Delta_t|$ at a time, when variables are
 branched.
 \item \textbf{Shortfall and energy on the reserve}, the pair $(s,E_W)$ of
 Proposition~\ref{prop:reserve} with the bound $2^{|W_0|}e^{-2s^2/E_W}$ and the
 exact transitions $s\mapsto s\pm\frac12b_w$, $E_W\mapsto E_W-b_w^2$ of
 \cite[Section~7]{CC}. This is the measure for designed cancellation.
 \item \textbf{Alignment.} The number of distinct signed neighbourhood types
 on the reserve, or the number of near-twin pairs; the measure of the
 structural exit, with cost $(n+1)^t$ for $t$ exact classes~\cite{N5} and
 $(3/4)^{m/(r+1)}$ for $m$ near-twin pairs~\cite{NT}.
\end{enumerate}

The first and the third are what Theorem~\ref{thm:biased} needs. The second
and the fourth are the two faces of the sum check, exact and Gaussian. The
fifth is what is left when the sum fails.

\section{Cleaning up: what to keep, and in which order}\label{sec:cleanup}

\paragraph{Core, to be read in this order.}
\begin{enumerate}
 \item \cite{N3}: the rules (domination, pruning clauses), the level budget,
 the tight bound for absolute linear gaps and the lower bound $C_{n,c}$.
 \item \cite{N4}: pair decisions, the sum rule and its characterization as
 the linear-relaxation cut, the clause weights, averaging over paths, the
 barrier for single-flip rules.
 \item \cite{N7,N8}: the race identity, Knuth's identity and the Key Lemma
 for random signs.
 \item \cite{EE,CC,CA}: early enumeration, the deterministic counting
 criterion and its checkable form.
 \item \cite{NT} and \cite[Section~5]{N5}: the structural exits.
 \item This note: the race for every rule, Theorem~\ref{thm:biased}, the
 three regimes, the reserve rule and the measures.
\end{enumerate}

\paragraph{Parked.} \cite{N6} entirely, except its $K_m$ example, which is in
\cite{N4}, and its joint refutation test, which is \cite[Proposition~7]{N4}.
\cite{N5} except Sections~3 and~5. \cite{N9}: its theorem is a special case of
the counting criterion of~\cite{CC,CA} with (C1) in spectral form, and the
biased part of its class is covered by Theorem~\ref{thm:biased} without any
norm.

\paragraph{One framework.} All kept results are statements about one search,
pruning clauses with the sum check (domination, propagation and pair decisions
as sound additions), counted by Knuth's identity and analysed by the race
$Z=S+Q-U$ along the uniform walk. The three regimes of Section~\ref{sec:regimes}
say which term wins. The open problem is Problem~\ref{prob:dichotomy}.

\paragraph{Next steps, in order.}
\begin{enumerate}
 \item Write Theorem~\ref{thm:biased} for $B=2$ and unit clauses in the
 input, which needs the race identity with multiplicities in written form
 (\cite[Section~4]{N9}, item~1), and optimize the constant with a single
 exponential moment of $Z_T$.
 \item Test Rule~R against the fixed order on balanced random formulas and on
 formulas with planted aligned groups, measuring $s^2/E_W$ at the reserve;
 the prediction is a constant-factor gain in the exponent on random formulas
 and a visible exit condition on planted ones.
 \item Prove Problem~\ref{prob:dichotomy} for the reserve: either Rule~R keeps
 $E_W\le C\,s\,\sqrt{|W_0|}$, or the rows $J_{x,W_0}$ of the remaining
 variables fall into few signed types. The first alternative is a vector
 balancing statement with chosen order and forced signs; the second is the
 inverse question of~\cite{NT}.
 \item Partial bias with the adaptive measure: recover the exponent
 $\max\{1-\eta,H((1-\varepsilon)/2)\}$ of~\cite[Theorem~11]{N3} in the relative
 setting, using Theorem~\ref{thm:biased} for the biased variables and the
 clause weights for the rest.
\end{enumerate}


\section{A plan for all formulas of growing density}\label{sec:plan}

The target is every formula with $m\le\rho n$, $\rho=\omega(1)$. Two facts fix
the shape of any such result. Formulas with $2^{\,n-O(n\log d/d)}$ strict local
maxima exist at every degree $d$ (disjoint copies of Max-Cut on $K_{d+1}$, and
Hastings' connected versions), and every rule implied by single flips keeps all
of them~\cite[Section~9]{N4}. So the algorithm needs two engines and the proof
needs a dichotomy: a \emph{race engine} (pruning clauses and the sum check,
which wins whenever the bad values cancel) and an \emph{aggregation engine}
(counting over groups of variables that act alike), and a lemma that every
node of every formula is handled by one of them. Theorem~\ref{thm:biased} and
the Key Lemma are the race engine on the two kinds of formulas where
cancellation is free or random. What follows is the plan for the rest.

\subsection{Cancellation in clause terms: the gap signs}

Along the uniform walk, with $g_s=\sgn\Delta_s$ the sign of the gap of $x_s$ at
its branching time and $\xi_s$ the fair coin of Lemma~\ref{lem:coins}(a),
\[
 b_w=\sum_{s\sim w,\,s\le T}\xi_s\,J_{sw}\,(-g_s)
 =-\tfrac12\sum_{s\sim w,\,s\le T}J_{sw}g_s\ +\ M_w ,
\]
where $M_w$ is a martingale with at most $m_w$ steps of size $\frac12$. So
$|b_w|$ is of order $\sqrt{m_w}$ exactly when the gap signs of the neighbours
of $w$, weighted by $J_{sw}$, cancel. The gap sign $g_s$ is the sign of
$h_s-\sum_{u<s}J_{su}a_u$, a function of the values of the earlier neighbours
of $x_s$; two gap signs $g_s,g_{s'}$ are correlated through the earlier
neighbours they share, and by a Littlewood--Offord argument
(condition on the shared coins; $|\E\sgn(Z+c)|\le\Pr[|Z|\le|c|]$ for the
private sum $Z$)
\[
 \bigl|\E[g_sg_{s'}]\bigr|\le C\,\frac{|\langle J_{s,\cdot},J_{s',\cdot}\rangle_{<}|+1}
 {\sqrt{D^-_sD^-_{s'}}},
\]
where $\langle J_{s,\cdot},J_{s',\cdot}\rangle_<$ is the \emph{signed codegree}
of $s$ and $s'$ over the earlier variables: common neighbours with agreeing
parities minus those with opposite parities. Summing over the pairs of
neighbours of $w$,
\[
 \E\Bigl(\sum_{s\sim w}J_{sw}g_s\Bigr)^2\ \lesssim\ m_w+\frac{C}{d}
 \sum_{s\ne s'\in N(w)}\bigl|\langle J_{s,\cdot},J_{s',\cdot}\rangle\bigr| ,
\]
and to leading order, with signs, the second term is the signed count of
$4$-cycles through $w$: pairs of clauses at $w$ and pairs of clauses at a
common neighbour $u$, counted with the parity of the cycle, which is the
diagonal entry $(J^4)_{ww}$. This is the combinatorial content of the spectral
condition: $\sum_w(J^4)_{ww}=\operatorname{tr}J^4\le\|J\|^2\,nd$, so
$\|J\|\le K\sqrt d$ bounds the average signed $4$-cycle count through a
variable by $K^2d^2$, against $d^3$ for $K_{d+1}$ and for $C_{n,c}$.

\begin{definition}[Alignment of a variable]
$\mathrm{al}(w)=\sum_{s\ne s'\in N(w)}J_{sw}J_{s'w}\langle J_{s,\cdot},J_{s',\cdot}\rangle$,
the number of even $4$-cycles through $w$ minus the number of odd ones; and
$\mathrm{al}^+(w)$ the same sum with absolute values.
\end{definition}

\paragraph{Target T1 (the race engine for fixed formulas).} Every formula with
$\mathrm{al}^+(w)\le\kappa d_wd$ for every $w$, or with
$\mathrm{al}(w)\le\kappa d_wd$ under a control of the lower-order terms, is
solved by the search of Section~\ref{sec:race} in time $2^{(1-\gamma(\kappa))n}$,
at every density. The proof is that of the Key Lemma~\cite{N8} with the
sign-correlation bound in place of the late signs, together with the
shortfall of Proposition~\ref{prop:gaps-balanced}. The class contains random
formulas, Max-Cut with arbitrary signs on every graph of codegree $O(1)$ (for
example girth at least $5$), and every formula with $\|J\|\le K\sqrt d$ on
average; its definition is a local count. The absolute version is enough for
$d\le\sqrt n$; for denser formulas the signed version is needed, and its
error terms are the first thing to work out.

\subsection{Aggregation when the gap signs are correlated}

If $\mathrm{al}^+(w)$ is large, then many pairs $s,s'$ of neighbours of $w$
have signed codegree of order $d$: similar signed neighbourhoods, hence
correlated gap signs, which is what defeats the sum. The two exits we have
need almost identical neighbourhoods: near-twin compression needs discrepancy
$r=O(1)$~\cite{NT}, exact classes need discrepancy $0$~\cite{N5}. The missing
tool is an aggregation for groups of variables with pairwise signed codegree
$(1-\kappa)d$: write $J=A+N$ with $A$ constant on the blocks of the groups and
$N$ the discrepancies, enumerate the group counts ($2^{O(n\log L/L)}$ choices
for groups of size $L$), which fixes the contribution of $A$, and search the
residual formula $N$, which is unaligned. The obstacle is soundness: pruning
clauses come from single flips, and a flip changes a group count. Inside a
count-constrained subproblem the pruning rule must therefore come from
\emph{swaps}, a variable and a group-mate flipped in opposite directions; this
is the class-flip rule of~\cite[Proposition~2]{N4}, whose window
(\cite[Proposition~1]{N4}) is the pruning clause of a swap, with the shared
neighbourhood cancelled and a cardinality condition on the discrepancies.

\paragraph{Target T2.} Unions of aligned blocks (cliques with constant $J$,
or copies of one signed neighbourhood) joined by a residual of degree $r$:
time $2^{o(n)}$ times the cost of the residual search. For $r=O(1)$ this is
\cite{NT}; for $r=\omega(1)$ it needs swap pruning on the residual with the
sum check, analysed by the race with the count constraints as live clauses
(they add to $m^0_w$ and $m^1_w$ equally, so they do not disturb $b_w$).

\subsection{The dichotomy and its first test}

\paragraph{Target T3.} If the average of $\mathrm{al}^+(w)$ over the free
variables exceeds $\kappa d_wd$ at a node, then a linear number of free
variables lie in groups of pairwise signed codegree at least $(1-\kappa')d$.
This is a statement about the graph of pairs with large signed codegree, of
K\H{o}v\'ari--S\'os--Tur\'an or dependent-random-choice type, and it is the
risky step: a blow-up in which each variable has $L$ copies with their
neighbourhoods re-signed at random on a $(1-\kappa)$-fraction of the positions
has pairwise codegree $\kappa d$ for $\kappa$ well below $1$, defeats the sum
for $L\to\infty$ (the aggregate on a free variable grows like $\sqrt{L m_w}$),
and is not compressible by near twins. The first experiment is therefore the
blow-up family with $L\in\{2,4,8,16\}$ and $\kappa\in\{0.5,0.8,0.95\}$:
measure $E_W=\sum_wb_w^2$ and the margin of the race, then the same with
count enumeration and swap pruning. Its outcome decides the form of T3.

\paragraph{Milestones.} (M1) T1 in the absolute form, $d\le\sqrt n$: the first
worst-case class of growing density defined by a local count and strictly
beyond random signs. (M2) T2 with swap pruning, first for bounded $r$ where it
must reproduce~\cite{NT}, then for $r=o(\sqrt d)$. (M3) The blow-up
experiment and T3 in the form it supports. (M4) The algorithm: at every node
compute $\mathrm{al}^+(w)$ for the free variables, run the race engine when it
is small and otherwise extract correlated groups and aggregate; the proof is
T1+T2+T3. If T3 fails, the fallback is a theorem for the union of the T1 class
and the near-twin class, both of unbounded density and both defined by
counts of clauses. The measures of Section~\ref{sec:measures} are the
quantities this plan monitors: $\mathrm{al}^+(w)$ decides the engine, $E_W$
tracks the race, and the group size $L$ prices the aggregation.

\section{Checks}\label{sec:checks}

\texttt{explore\_biased\_class.py} reuses the search and the population
splitting estimator of \texttt{explore\_key\_lemma.py}.
\begin{itemize}
 \item Without arguments (minutes): Table~\ref{tab:biased}, with the
 verification that every generated formula satisfies (B1) for the stated
 $\varepsilon$ (the minimum ratio $|h_v|/D_v$ is printed).
 \item \texttt{terms} (a minute): the race terms at depth $0.8n$ and $0.9n$
 on biased formulas with $n=1000$, $d\in\{6,30,200,500\}$ and
 $\varepsilon\in\{\frac14,\frac12\}$, in units of $n\bar d$: $S_T$ between
 $0.08$ and $0.30$, $U_T$ between $0.004$ and $0.04$, $Q_T$ within $\pm0.002$.
 The race is won by a factor $5$ to $20$ at every density, which is the
 mechanism of the proof; the proof's margin $\frac2{32}\varepsilon\eta$ is far
 below it.
\end{itemize}

\begin{thebibliography}{99}
\raggedright
\bibitem{N2} Exact Max 2-SAT with an absolute linear polarity gap. Working
 note, 23 September 2026 (\texttt{max2sat\_absolute\_linear\_gap\_proof.tex}).
\bibitem{N3} Max 2-SAT and the domination rule: results and ideas so far.
 Working note, 5 October 2026 (\texttt{max2sat\_results\_and\_ideas.tex}).
\bibitem{N4} Max 2-SAT with pruning clauses: pair decisions, sums, and a
 measure. Working note, 6--7 October 2026
 (\texttt{max2sat\_pairs\_sums\_measure.tex}).
\bibitem{N5} Exploring the branching frontier for Max 2-SAT. Draft,
 7 October 2026 (\texttt{max2sat\_branching\_frontier.tex}).
\bibitem{N6} A learning and coverage dichotomy for Max 2-SAT. Draft,
 7 October 2026 (\texttt{max2sat\_learning\_dichotomy.tex}).
\bibitem{N7} A simplified plan for Max 2-SAT: one search, one count, one Key
 Lemma. Working note, 8 October 2026 (\texttt{max2sat\_plan\_key\_lemma.tex}).
\bibitem{N8} The Key Lemma for random signs. Working note, 8 October 2026
 (\texttt{max2sat\_key\_lemma\_proof.tex}).
\bibitem{N9} Max 2-SAT below $2^n$ on unaligned formulas. Working note,
 8 October 2026 (\texttt{max2sat\_explicit\_class.tex}).
\bibitem{EE} Max 2-SAT by early constraint-guided enumeration. Working note,
 8 October 2026 (\texttt{max2sat\_early\_enumeration.tex}).
\bibitem{CC} Replacing random signs by neighborhood counts. Working note,
 8 October 2026 (\texttt{max2sat\_counting\_conditions.tex}).
\bibitem{CA} A checkable replacement for the gap-profile condition. Addendum,
 8 October 2026 (\texttt{max2sat\_counting\_addendum.tex}).
\bibitem{NT} Near-twin neighborhoods as a structural exit. Working note,
 8 October 2026 (\texttt{max2sat\_near\_twin\_pairs.tex}).
\bibitem{Azu67} K.~Azuma. Weighted sums of certain dependent random variables.
 \emph{T\^ohoku Mathematical Journal} 19(3):357--367, 1967.
\bibitem{Hoe63} W.~Hoeffding. Probability inequalities for sums of bounded
 random variables. \emph{Journal of the American Statistical Association}
 58(301):13--30, 1963.
\bibitem{Mau03} A.~Maurer. A bound on the deviation probability for sums of
 non-negative random variables. \emph{Journal of Inequalities in Pure and
 Applied Mathematics} 4(1), Article~15, 2003.
\end{thebibliography}

\end{document}
